\scanpage{312}
- \mspage{20} -
\margin{n° \unsure{2|1}69}

les de $G'$ et $G'_1$) et que $u_1=f\circ u$.

2) Lorsque $A=B=G$ (la structure de $(A,A)$-bimodule de $A$ étant définie par les homothéties à gauche et à droite), et que $A'=B'$ et $h=h'$ on peut prendre pour $G'$ l'anneau $A'$ et pour $u$ l'homomorphisme $h$ de $A$ dans $A'$.

3) Supposons que l'on ait $A=B$, $A'=B'$, $h=h'$, que les anneaux $A$ et $A'$ soient \uline{commutatifs}, et que la structure de $A$-module à gauche de $G$ coincide avec sa structure de $A$-module à droite. On peut alors prendre pour $G'$ le produit tensoriel $A'\otimes_A G$ (la structure de $A'$-module à droite de $G'$ coincidant avec sa structure de $A'$-module à gauche) et pour $u$ l'application $g\to1\otimes g$ ($g\in G$) de $G$ dans $G'$. Nous dirons alors que la forme bilinéaire (resp. sesquilinéaire) $\Phi'$ définie par la prop. 1 (resp. prop.2) est obtenue à partir de $\Phi$ \uline{par extension de l'anneau de base}, ou \uline{par extension des scalaires}.

Ce qui suit est valable aussi bien dans le cas bilinéaire que dans le cas sesquilinéaire ; les hypothèses et notations sont celles de la prop.1 (resp. prop.2). Etant donné un sous module $M$ de $E$ ou $F$, on désignera par $M'$ le sous module de $E'$ ou $F'$ engendré par l'image canonique de $M$.

PROPOSITION 3 - \uline{Les hypothèses et notations étant celles de la prop.1} (resp. \uline{prop.2}) \uline{supposons de plus que} $A,B,A',B'$ \uline{soient des corps et que les applications} $\alpha$ \uline{et} $\beta$ \uline{de} $A'\otimes_A \uline{G}$ \uline{et} $G\otimes_A B'$ \uline{dans} $G'$ \uline{caractérisées par} $\alpha(a'\otimes g)=a'u(g)$ \uline{et} $\beta(g\otimes b')=u(g)b'$ ($a'\in A'$, $b'\in B'$, $g\in G$) \uline{soient injectives}. \uline{Soient} $M$ \uline{un sous}
\mspage{311}
